% Part: computability
% Chapter: recursive-functions
% Section: primes

\documentclass[../../../include/open-logic-section]{subfiles}

\begin{document}

\olfileid{cmp}{rec}{pri}
\olsection{الأعداد الأولية}

بالتكميم المحدود والتصغير المحدود يتيسر لنا إثبات العودية البدائية لكثير من الدوال
والعلائق المألوفة. فلننظر، مثلًا، في العلاقة «$x$ يقسم
$y$»، التي تكتب $x\mid y$. ومعنى $x\mid y$ أن تقبل قسمة
$y$ على~$x$ دون باق، أي أن يكون $y$ مضاعفًا صحيحًا لـ~$x$. (وإذا لم
تصدق، أي إذا كان باقي قسمة $y$ على $x$ أكبر من~$0$، كتبنا
$x\nmid y$.) وبعبارة أخرى، $x\mid y$ إذا وفقط إذا وجد~$z$ بحيث
$x\cdot z=y$. ومتى وجد شاهد $z$ لهذا، أمكن اختيار شاهد يكون
$\leq y$؛ حتى في حالة كون العددين صفرًا نختار الصفر شاهدًا. ومن ثم يصدق $x\mid y$ إذا وفقط إذا وجد $z\le y$ بحيث
$x\cdot z=y$. فتعرّف العلاقة $x\mid y$ بالتكميم الوجودي
المحدود من المساواة والضرب على الوجه الآتي:
\[
x \mid y \defiff \bexists{z \leq y}{(x \cdot z) = y}.
\]
فتثبت بذلك العودية البدائية لـ$x\mid y$.

العدد الطبيعي~$x$ \emph{أولي} متى كان غير $0$ وغير $1$، ولم
يقبل القسمة إلا على $1$ ونفسه. فشرط القواسم في العدد
الأولي أن $y\mid x$ يستلزم إما $y=1$ وإما~$y=x$. فإذا استبعدنا الصفر والواحد وأردنا
اختبار أولية $x$، كفانا فحص $y\mid x$ لكل $y\le x$، لأن
$y\nmid x$ تلقائيًا إذا كان $y>x$. وعليه يكون تعريف العلاقة
$\fn{Prime}(x)$، التي تصدق إذا وفقط إذا كان $x$ أوليًا، بـ
\[
\fn{Prime}(x) \defiff x \geq 2 \land \bforall{y \leq x}{(y \mid x \lif y
  = 1 \lor y = x)}
\]
فتكون عودية بدائية.

الأعداد الأولية هي $2$ و$3$ و$5$ و$7$ و$11$، وهكذا. ولنسمّ الدالة
$p(x)$ التي تعطي العدد الأولي ذا الرتبة $x$ في هذه المتتالية؛ أي
$p(0)=2$ و$p(1)=3$ و$p(2)=5$، وهكذا. (وتيسيرًا للكتابة سنكتب غالبًا
$p_x$ بدلًا من $p(x)$؛ فـ$p_0=2$ و$p_1=3$، وهكذا.)

فإن حصلت لنا دالة $\fn{nextPrime(x)}$ تعطي أول عدد أولي يجاوز~$x$،
سهل تعريف $p$ بالعودية البدائية:
\begin{align*}
  p(0) & = 2\\
  p(x+1) & = \fn{nextPrime}(p(x))
\end{align*}
وقيمة $\fn{nextPrime}(x)$ هي أصغر $y$ يجمع الشرطين: $y>x$ وكون $y$ أوليًا؛ فلا يعسر
طلبها بالبحث غير المحدود. لكننا نستغني عن ذلك، ونعرفها بالتصغير
المحدود، اعتمادًا على نتيجة لإقليدس: بين $x$ و
$\fact{x}+1$ دائمًا عدد أولي.
\[
  \fn{nextPrime(x)} =
  \bmin{y \leq \fact{x}+1}{(y > x \land \fn{Prime}(y))}.
\]
فتكون $\fn{nextPrime}(x)$، ويتبعها $p(x)$، عوديتين بدائيتين،
وهو أقوى من مجرد قابليتهما للحساب.

(وبرهان نتيجة إقليدس موجز: إذا كان $x<2$، فالعدد الأولي~$2$ أكبر
من~$x$ ولا يجاوز $\fact{x}+1$. وأما إذا كان $x\geq2$، فليكن $p_n$
أكبر عدد أولي $\le x$، ولنأخذ الجداء
$p=p_0\cdot p_1\cdot\dots\cdot p_n$ لجميع الأعداد الأولية~$\le x$.
إما أن يكون $p+1$ أوليًا، وإما أن يوجد عدد أولي بين $x$ و~$p+1$.
وبيانه أنه إن لم يكن $p+1$ أوليًا، فلا بد من عدد أولي $q\mid p+1$
بحيث $q<p+1$. ولا يقسم أي عدد أولي $\le x$ العدد $p+1$. (فبحسب
تعريف~$p$، يقسم كل عدد أولي $p_i\le x$ العدد~$p$ بلا باق. ولذلك تكون
قسمة $p+1$ على كل $p_i\le x$ ذات باق~$1$، ومن ثم
$p_i\nmid p+1$.) وهكذا يكون $q$ عددًا أوليًا أكبر من~$x$ وأصغر من
$p+1$. ولأن $p\le\fact{x}$، يوجد عدد أولي أكبر من~$x$ وأصغر من أو
مساو لـ$\fact{x}+1$.)

% تصحيح مصدر (C219-01): فُصلت حالتا x=0 وx=1 قبل اختيار أكبر أولي لا يجاوز x.
\textbf{تنبيه تصحيحي على الأصل.} فُصلت حالتا $x=0$ و$x=1$؛ إذ لا
يتأتى فيهما اختيار «أكبر عدد أولي لا يجاوز~$x$»، ويشهد لهما العدد~$2$
مباشرة.

\begin{prob}
ضع تعريفًا للقسمة الصحيحة $d(x,y)$ بالتصغير المحدود.
\end{prob}

\end{document}
